\chapter{Planificación}
\label{cap:planificacion}

Este capítulo recoge la planificación seguida para desarrollar el trabajo
entre el 20 de julio y el 23 de agosto de 2026 (cinco semanas), organizada
en siete fases secuenciales, desde la selección de la metodología de
referencia y del entorno de pruebas hasta la redacción final de la memoria.
Al tratarse de un trabajo individual, las fases se han ejecutado de forma
mayormente secuencial, con solapamientos puntuales entre fases contiguas
(por ejemplo, entre el diseño del entorno y la definición de la metodología
propia, desarrollados con fuerte dependencia mutua).

\section{Fases de trabajo}

\begin{enumerate}
\item \textbf{Selección de la metodología de referencia y del entorno de
  pruebas.} Estudio preliminar de metodologías de seguridad ofensiva
  existentes (OWASP WSTG, PTES, OSSTMM, NIST SP 800-115) y elección de la
  OWASP WSTG como referencia; decisión de construir un entorno de pruebas
  propio (\breachlab) en lugar de partir de uno ya existente, para poder
  instrumentar el código y comparar ambas ramas, vulnerable y corregida.
\item \textbf{Estudio de la metodología de referencia.} Revisión en
  profundidad de la estructura, fases y catálogo de pruebas de la OWASP
  WSTG (capítulo~\ref{cap:marco-teorico}).
\item \textbf{Diseño e implementación del entorno de pruebas.} Definición de
  la arquitectura de \breachlab, implementación del backend y el frontend,
  introducción deliberada de las seis vulnerabilidades y de su
  contramedida, y despliegue mediante Docker Compose
  (capítulo~\ref{cap:diseno}).
\item \textbf{Definición de la metodología propia.} Derivación de una
  metodología ágil a partir de la OWASP WSTG, acotando su alcance a un
  subconjunto de pruebas representativo (capítulo~\ref{cap:metodologia}).
\item \textbf{Ejecución de las pruebas.} Aplicación de la metodología propia
  sobre \breachlab: identificación, explotación, evidencia y verificación
  de la contramedida de cada vulnerabilidad (capítulo~\ref{cap:explotacion});
  comparación con dos laboratorios de referencia y con la propia OWASP WSTG
  (capítulo~\ref{cap:comparativa}); y construcción de un caso de estudio que
  encadena varias vulnerabilidades (capítulo~\ref{cap:caso-estudio}).
\item \textbf{Análisis de resultados y redacción de conclusiones.} Valoración
  del grado de cumplimiento de los objetivos planteados en el
  apartado~\ref{sec:objetivos} y propuesta de líneas de trabajo futuro
  (capítulo~\ref{cap:conclusiones}).
\item \textbf{Redacción y revisión final de la memoria.} Consolidación de
  todos los capítulos, revisión de forma y contenido, y preparación de los
  anexos de instalación y manual de usuario.
\end{enumerate}

\section{Estimación temporal}

La figura~\ref{fig:gantt-planificacion} muestra el diagrama de Gantt de las
siete fases descritas, a lo largo de las cinco semanas del proyecto, y la
tabla~\ref{tab:planificacion} recoge la duración de cada fase y sus fechas
concretas.

\begin{figure}[htbp]
  \centering
  \resizebox{\textwidth}{!}{%
  \begin{ganttchart}[
      hgrid,
      vgrid={*{7}{draw=none}, draw=gray!40},
      x unit=0.38cm,
      y unit title=0.7cm,
      y unit chart=0.6cm,
      title height=1,
      title/.style={draw=none, fill=gray!25},
      title label font=\scriptsize\bfseries,
      bar height=0.55,
      bar/.style={fill=gray!70},
      bar label font=\scriptsize,
      bar label node/.append style={text width=4.4cm, align=right},
    ]{1}{35}
    \gantttitle{Semana 1}{7}
    \gantttitle{Semana 2}{7}
    \gantttitle{Semana 3}{7}
    \gantttitle{Semana 4}{7}
    \gantttitle{Semana 5}{7} \\
    \ganttbar{1. Selección de metodología y entorno}{1}{2} \\
    \ganttbar{2. Estudio de la metodología de referencia}{3}{6} \\
    \ganttbar{3. Diseño e implementación de \breachlab}{7}{15} \\
    \ganttbar{4. Definición de la metodología propia}{16}{19} \\
    \ganttbar{5. Ejecución de pruebas y comparativa}{20}{28} \\
    \ganttbar{6. Análisis de resultados y conclusiones}{29}{30} \\
    \ganttbar{7. Redacción y revisión final}{31}{35}
  \end{ganttchart}%
  }
  \caption{Diagrama de Gantt de la planificación.}
  \label{fig:gantt-planificacion}
\end{figure}

\begin{table}[htbp]
  \centering
  \small
  \begin{tabular}{p{6.4cm}p{1.6cm}p{2.5cm}p{2.5cm}}
    \toprule
    \tabhead{Fase} & \tabhead{Duración} & \tabhead{Fechas} & \tabhead{Entregable} \\
    \midrule
    1. Selección de metodología y entorno & 2 días & 20--21 jul. & Alcance definido \\
    2. Estudio de la metodología de referencia & 4 días & 22--25 jul. & Capítulo~\ref{cap:marco-teorico} \\
    3. Diseño e implementación de \breachlab & 9 días & 26 jul.--3 ago. & \breachlab desplegado \\
    4. Definición de la metodología propia & 4 días & 4--7 ago. & Capítulo~\ref{cap:metodologia} \\
    5. Ejecución de pruebas y comparativa & 9 días & 8--16 ago. & Capítulos~\ref{cap:explotacion}--\ref{cap:caso-estudio} \\
    6. Análisis de resultados y conclusiones & 2 días & 17--18 ago. & Capítulo~\ref{cap:conclusiones} \\
    7. Redacción y revisión final & 5 días & 19--23 ago. & Memoria final \\
    \bottomrule
  \end{tabular}
  \caption{Planificación temporal del trabajo.}
  \label{tab:planificacion}
\end{table}

Las fases de mayor carga son la 3 (diseño e implementación de \breachlab) y
la 5 (ejecución de las pruebas y comparativa), al concentrar la parte más
técnica del trabajo; el resto de fases, de estudio y de redacción, se han
planificado con una duración más ajustada. Las fases 3 y 4 se han
desarrollado además con un solapamiento parcial: varias decisiones de
diseño de \breachlab (por ejemplo, qué vulnerabilidades incluir) se han
tomado ya pensando en el subconjunto de pruebas que definiría la
metodología propia, por lo que ambas fases se han retroalimentado en lugar
de ejecutarse de forma estrictamente independiente.
